\section{Discussion, limitations and reflection}
\subsection*{Theory versus practice}
NP-hardness concerns exact answers on every instance. It does not prevent finding good plans
quickly on practical inputs. TS often reaches the reference early in these runs, although it
does not provide an optimality certificate for the team instance.

Approximation methods offer a proved worst-case relationship to the optimum under stated
assumptions. Tour partitioning starts from a TSP tour and cuts it into capacity-feasible routes
[8]. The starting tour's quality affects the guarantee. Fleet size needs separate attention:
capacity-feasible cuts can still use too many vehicles. Metaheuristics have no general quality
guarantee, but can search with the project's constraints and a fixed runtime. In practice,
measured quality, repeatability and independent feasibility checks complement theoretical guarantees.

\subsection*{Limitations}
Five core instances on one machine do not establish universal TS superiority. Demands are
deterministic. The model omits battery range, wind, flight permissions and delivery time windows.
It is an algorithm study. The team reference is empirical, and shared hardware affects timing.

\subsection*{Reflection grounded in the project records}
The clearest surprise is that the crossing penalty does not reorganize the synthetic instance.
Its saved plans already achieve the minimum crossing-route count required by European demand.

Earlier implementation notes record a GA replacement-order bug: after moving a replacement
left in a sorted population, the second ordering loop restarted from its old position. This
could break the sorted-population invariant and weaken elitism. The recorded fix continues
from the updated position in \texttt{ga.cpp}. Maintaining this representation invariant matters
even when a run's final cost looks plausible.

A useful next step is a controlled comparison with a small local search inside GA, using the
existing memetic variant as a starting point. Battery range or time windows could follow after
the core algorithms are fully understood.

\subsection*{References}
\begingroup\small\setlength{\parskip}{3pt}
[1] Garey, M. R. and Johnson, D. S. (1979). \emph{Computers and Intractability}. Freeman.

[2] Clarke, G. and Wright, J. W. (1964). Scheduling of vehicles from a central depot to a
number of delivery points. \emph{Operations Research}, 12(4), 568--581.

[3] Prins, C. (2004). A simple and effective evolutionary algorithm for the vehicle routing
problem. \emph{Computers \& Operations Research}, 31(12), 1985--2002.

[4] Vidal, T. (2016). Technical note: Split algorithm in $O(n)$ for the capacitated vehicle
routing problem. \emph{Computers \& Operations Research}, 69, 40--47.

[5] CVRPLIB. Capacitated Vehicle Routing Problem Library. \url{http://vrp.atd-lab.inf.puc-rio.br/}.

[6] ML4CO. \emph{VRPLIBOriDataset}, cached benchmark mirror.
\url{https://huggingface.co/datasets/ML4CO/VRPLIBOriDataset}.

[7] Glover, F. and Laguna, M. (1997). \emph{Tabu Search}. Kluwer.

[8] Haimovich, M. and Rinnooy Kan, A. H. G. (1985). Bounds and heuristics for capacitated
routing problems. \emph{Mathematics of Operations Research}, 10(4), 527--542.

[9] Davis, L. (1985). Applying adaptive algorithms to epistatic domains. \emph{IJCAI}, 162--164.
\endgroup
